\documentclass[10pt,a4paper]{amsart}
\usepackage[margin=19mm]{geometry}
\usepackage{amsmath,amssymb,mathtools}
\usepackage[scr=rsfs,cal=euler]{mathalfa}
\usepackage{xfrac}
\usepackage[unicode,hidelinks,bookmarks=false]{hyperref}
\newcommand{\ie}{i.e.\ }
\newcommand{\cf}{cf.\ }
\newcommand{\resp}{respectively\ }
\newcommand{\Subsection}{\subsection}
\providecommand{\nobreakdash}{\nobreak\hbox{-}}
\setlength{\emergencystretch}{2em}
\multlinegap0pt
\usepackage{algorithm}\usepackage{algpseudocode}
\usepackage{amssymb}
\usepackage{bm}
\usepackage{float}
\newtheorem{lem}{Lemma}
\title{Some Tur\'{a}n-type results for {the} signless Laplacian spectral radius}
\author{Jian Zheng, Yongtao Li, Yi-Zheng Fan}
\date{Source: arXiv:2507.02263v2 (CC BY 4.0)}
\begin{document}
\maketitle

\noindent\emph{Source and licence.} Some Tur\'{a}n-type results for {the} signless Laplacian spectral radius. Version arXiv:2507.02263v2 (CC BY 4.0), distributed by arXiv under the Creative Commons Attribution 4.0 International licence, http://creativecommons.org/licenses/by/4.0/, as recorded in the arXiv licence metadata for this exact version and shown on the abstract page https://arxiv.org/abs/2507.02263v2. Full licence notice: \texttt{licenses/arxiv-2507.02263-CC-BY-4.0.txt}.\par\smallskip

\noindent\textit{Editorial scope.} Counted unit: the article's lem gj (source0.tex lines 1376--1380). The article's complete original proof of it is delivered with it (source0.tex lines 1382--1412, one \texttt{proof} environment of 31 lines). No mathematics is written, repaired or re-derived: the statement and the proof are the article's own lines, in the article's own notation, and no retained line is substituted or re-flowed.  No further source-local context is needed: every cross-reference of every retained line resolves inside the two delivered spans.  Nothing was omitted from any retained span, so the package carries no omission record.  No retained line contains a citation, so the wrapper carries no bibliography and no bibliography entry.  No retained line contains a figure environment, an image-inclusion command or drawing code, so no image is delivered and the package declares no asset dependency.  Editorial formatting only, in addition: an AMS \texttt{amsart} preamble that restates the article's own declarations and macros used by the retained lines, namely the environments \texttt{lem} on separate counters, together with the standard packages those lines need.  The retained environments print in this document as Lemma 1 in order of appearance, whereas the article prints the same items with the numbering of its own full document.  The complete native source of the article as distributed in the arXiv e-print for this exact version is bundled unchanged as \texttt{sources/180935/source0.tex}.
\begin{lem}\label{gj}
Let $H$ be a linear $r$-graph. Then
$$q(H)\leq\frac{1}{\binom{r-1}{s-1}}q(\partial_{s} H).$$
In particular, we have $(r-1)q(H)\leq q(\partial H)$.
\end{lem}

\begin{proof}
Let $\mathbf{x}$ be a nonnegative eigenvector corresponding to $q(H)$. 
Then  $\sum_{i\in V(H)} x_i^r =1$ and 
\[  q(H)=\sum_{i\in V(H)}d_{H}(i)x_{i}^{r}+r\sum_{e\in E(H)}x^{e}. \]
For each fixed $i\in V(H)$ and $e\in E(H)$ with $i\in e$, 
there are ${r-1 \choose s-1}$ subsets of $e$ with size $s$ and containing $i$. 
Then $d_H(i)= \frac{1}{{r-1 \choose s-1}}d_{\partial_s H}(i)$. Therefore, we have 
\begin{align}\label{gj1}
\begin{split}
q(H) = \frac{1}{\binom{r-1}{s-1}}\sum_{i\in V(\partial_{s}H)}d_{\partial_{s}H}(i)(x_{i}^{r/s})^{s}
+r\sum_{e\in E(H)}\prod_{\{i_{1},\ldots,i_{s}\}\subset e} 
(x_{i_{1}}\cdots x_{i_{s}})^{\frac{1}{\binom{r-1}{s-1}}}.
 \end{split}
\end{align} 
Using the AM-GM inequality and the Power-Mean inequality, we obtain 
\begin{align}
r\sum_{e\in E(H)}\prod_{\{i_{1},\ldots,i_{s}\}\subset e}(x_{i_{1}}\cdots x_{i_{s}})^{\frac{1}{\binom{r-1}{s-1}}}
&\leq  r\sum_{e\in E(H)}\bigg(\frac{1}{\binom{r}{s}}\sum_{\{i_{1},\ldots,i_{s}\}\subset e}(x_{i_{1}}\cdots x_{i_{s}})^{\frac{1}{\binom{r-1}{s-1}}}\bigg)^{\binom{r}{s}} \notag \\
&\leq  r\sum_{e\in E(H)} \frac{1}{\binom{r}{s}}\sum_{\{i_{1},\ldots,i_{s}\}\subset e}(x_{i_{1}}\cdots x_{i_{s}})^{\frac{\binom{r}{s}}{\binom{r-1}{s-1}}}  \notag \\
&= \frac{s}{\binom{r-1}{s-1}}\sum_{e\in E(H)} \sum_{\{i_{1},\ldots,i_{s}\}\subset e}(x_{i_{1}})^{\frac{r}{s}}\cdots(x_{i_{s}})^{\frac{r}{s}} \notag \\
&= \frac{s}{\binom{r-1}{s-1}} \sum_{\{i_{1},\ldots,i_{s}\}\in E(\partial_{s} H)}(x_{i_{1}})^{\frac{r}{s}}\cdots(x_{i_{s}})^{\frac{r}{s}}. \label{gj2}
\end{align}
Putting (\ref{gj1}) and (\ref{gj2}) together, we find that
\begin{align*} \notag
q(H)&\leq \frac{1}{\binom{r-1}{s-1}}\bigg(\sum_{i\in V(\partial_{s}H)}d_{\partial_{s}H}(i)(x_{i}^{r/s})^{s}+
s\sum_{\{i_{1},\ldots,i_{s}\}\in E(\partial_{s} H)}(x_{i_{1}})^{\frac{r}{s}}\cdots(x_{i_{s}})^{\frac{r}{s}}\bigg)\\ \notag 
&\leq \frac{1}{\binom{r-1}{s-1}}q(\partial_{s} H),
\end{align*}
where the last inequality follows from $\sum_{i\in V(\partial_{s} H)}(x_{i}^{r/s})^{s}=1$. 
This completes the proof. 
\end{proof}

\end{document}
